\lecture{4}{Graph Neural Networks II: A General Perspective}{Thursday, October 1, 2026}\label{ch:lec04}

% Local TikZ styles (prefix lec04)
\tikzset{
  lec04A/.style={fill=Goldenrod!70},
  lec04B/.style={fill=BrickRed!60},
  lec04C/.style={fill=ForestGreen!50},
  lec04D/.style={fill=RoyalBlue!45},
  lec04E/.style={fill=Orchid!60},
  lec04F/.style={fill=WildStrawberry!45},
  lec04n/.style={circle, draw=black!60, thick, minimum size=4.6mm,
                 inner sep=0pt, font=\scriptsize\sffamily},
  lec04agg/.style={draw=black!60, fill=black!10, minimum size=6mm},
  lec04agg1/.style={draw=black!70, fill=black!55, minimum size=4mm},
  lec04msg/.style={draw=black!50, fill=black!25, minimum width=2.8mm,
                   minimum height=1.1mm, inner sep=0pt},
  lec04in/.style={black!55, densely dotted, -{Stealth[length=1.3mm]}},
  lec04lay/.style={draw, thick, fill=white, minimum width=19mm,
                   minimum height=4.6mm, inner sep=1pt, font=\scriptsize\sffamily},
  lec04lbl/.style={font=\scriptsize\sffamily\bfseries},
}

\textbf{Outline.} Recap $\to$ a general GNN framework (5 design dimensions)
$\to$ a single GNN layer (message + aggregation; GCN, GraphSAGE, GAT) $\to$ GNN
layers in practice (BatchNorm, dropout) $\to$ stacking GNN layers
(over-smoothing, expressive shallow GNNs, skip connections).

% ---------------------------------------------------------------------
\section{Recap and a general GNN framework}\label{sec:lec04-framework}

\subsection{Recap: deep graph encoders and GNNs}\label{sec:lec04-recap}
\slides{3--6}

(Details in \cref{ch:lec03}.)
\begin{itemize}
  \item \textbf{Deep graph encoder:} graph convolutions $\to$ activation
    function $\to$ regularization (e.g.\ dropout) $\to$ graph convolutions
    $\to\cdots$ \textbf{Output:} node embeddings; we can also embed subgraphs,
    graphs.
  \item \textbf{Idea:} a node's neighborhood defines a \term{computation graph}
    ($k=1$, $k=2$ hops around node $i$): (1) determine node computation graph,
    (2) propagate and transform information (aggregator$_1$, aggregator$_2$ at
    the two levels).
  \item \textbf{Intuition:} nodes aggregate information from their neighbors
    using neural networks (every box of the computation graph is a NN).
  \item \textbf{Every node defines a computation graph based on its
    neighborhood!} (one tree per node $A,\dots,F$ of the input graph).
\end{itemize}
\begin{keyidea}
  Learn how to propagate information across the graph to compute node features.
\end{keyidea}

\subsection{A general GNN framework}\label{sec:lec04-design}
\slides{7--12}

Five parts of a GNN (You, Ying, Leskovec, \emph{Design Space of Graph Neural
Networks}, NeurIPS 2020), \cref{fig:lec04-framework}:
\begin{itemize}
  \item \textbf{(1) Message + (2) Aggregation:} \textbf{GNN layer = Message +
    Aggregation}. Different instantiations under this perspective: GCN,
    GraphSAGE, GAT, \dots\ (\cref{sec:lec04-layer}).
  \item \textbf{(3) Layer connectivity:} connect GNN layers into a GNN
    $\Rightarrow$ stack layers sequentially; ways of adding skip connections
    (\cref{sec:lec04-stack}).
  \item \textbf{(4) Graph augmentation:} idea: \textbf{raw input graph $\neq$
    computational graph} $\Rightarrow$ graph feature augmentation, graph
    structure augmentation.
  \item \textbf{(5) Learning objective:} how do we train a GNN?
    Supervised/unsupervised objectives; node/edge/graph-level objectives (all
    discussed later in the class).
\end{itemize}

\begin{figure}[htbp]
  \centering
  \begin{tikzpicture}[scale=0.85]
    % dashed layer boxes
    \draw[densely dotted, thick, BrickRed] (-3.1,2.4) rectangle (3.1,3.62);
    \draw[densely dotted, thick, BrickRed] (-3.1,0.1) rectangle (3.1,1.42);
    % computation graph of A
    \node[lec04n, lec04A] (r) at (0,4.15) {A};
    \node[lec04agg] (g2) at (0,3.2) {};
    \node[lec04n, lec04B] (B) at (-2.2,2.0) {B};
    \node[lec04n, lec04C] (C) at (0,2.0) {C};
    \node[lec04n, lec04D] (D) at (2.0,2.0) {D};
    \foreach \n in {B,C,D} {
      \node[lec04agg1] (g\n) at ($(\n)+(0,-1)$) {};
      \draw[-{Stealth[length=1.3mm]}, black!70] (g\n) -- (\n);
      \draw[lec04in] (\n) -- (g2) node[pos=0.5, lec04msg, sloped] {};
    }
    \draw[-{Stealth[length=1.3mm]}, black!70] (g2) -- (r);
    \foreach \c/\x/\p [count=\i] in {A/-2.7/B, C/-1.7/B, A/-1.05/C, B/-0.35/C,
                                     E/0.35/C, F/1.05/C, A/2.0/D} {
      \node[lec04n, lec04\c] (leaf\i) at (\x,-0.4) {\c};
      \draw[lec04in] (leaf\i) -- (g\p) node[pos=0.5, lec04msg, sloped] {};
    }
    % labels
    \node[lec04lbl, BrickRed, anchor=east] at (-3.2,3.0) {GNN Layer 2};
    \node[lec04lbl, BrickRed, anchor=east] at (-3.2,0.8) {GNN Layer 1};
    \node[lec04lbl, BrickRed, anchor=east, align=right] at (-3.2,1.92)
      {(3) Layer\\connectivity};
    \node[lec04lbl, black!65, anchor=west] at (3.2,3.25) {(2) Aggregation};
    \node[lec04lbl, black!65, anchor=west] at (3.2,2.6) {(1) Message};
    \node[lec04lbl, SkyBlue!70!black] at (0,4.75) {(5) Learning objective};
    \node[lec04lbl, RoyalBlue] at (0,-0.98) {(4) Graph augmentation};
    % input graph
    \begin{scope}[shift={(6.6,0.7)}, scale=0.75]
      \node[lec04n, lec04A] (a) at (0.9,1.25) {A};
      \node[lec04n, lec04B] (b) at (2.5,1.9) {B};
      \node[lec04n, lec04C] (c) at (2.9,1.0) {C};
      \node[lec04n, lec04D] (d) at (0.0,0.0) {D};
      \node[lec04n, lec04E] (e) at (1.9,0.1) {E};
      \node[lec04n, lec04F] (f) at (3.6,0.35) {F};
      \draw[edge] (a)--(b) (a)--(c) (a)--(d) (b)--(c) (c)--(e) (c)--(f) (e)--(f);
      \draw[arrow] (0.9,2.45) node[above, font=\scriptsize\sffamily] {target node}
        -- (0.9,1.6);
      \node[font=\scriptsize\sffamily] at (1.8,-0.75) {input graph};
    \end{scope}
  \end{tikzpicture}
  \caption{The general GNN framework on the 2-layer computation graph of target
  node $A$ (slides~7, 12).}
  \label{fig:lec04-framework}
\end{figure}

% ---------------------------------------------------------------------
\section{A single layer of a GNN}\label{sec:lec04-layer}

\subsection{Message and aggregation}\label{sec:lec04-msgagg}
\slides{14--19}

\begin{itemize}
  \item \textbf{Idea of a GNN layer:} compress a set of vectors into a single
    vector. \textbf{Two-step process:} (1) message, (2) aggregation.
  \item Input node embeddings $\vh_v^{(l-1)}, \vh_{u\in N(v)}^{(l-1)}$ (from node
    itself + neighboring nodes) $\to$ $l$-th GNN layer $\to$ output node
    embedding $\vh_v^{(l)}$.
  \item \textbf{(1) Message computation} -- \term{message function}; intuition:
    each node will create a message, which will be sent to other nodes later:
\end{itemize}
\begin{equation}\label{eq:lec04-msg}
  \vm_u^{(l)} = \MSG^{(l)}\big(\vh_u^{(l-1)}\big)
\end{equation}
\begin{itemize}
  \item[] \textbf{Example:} a linear layer $\vm_u^{(l)}=\mW^{(l)}\vh_u^{(l-1)}$
    (multiply node features with weight matrix $\mW^{(l)}$).
  \item \textbf{(2) Aggregation} -- intuition: node $v$ will aggregate the
    messages from its neighbors $u$:
\end{itemize}
\begin{equation}\label{eq:lec04-agg}
  \vh_v^{(l)} = \AGG^{(l)}\big(\{\vm_u^{(l)}, u\in N(v)\}\big)
\end{equation}
\begin{itemize}
  \item[] \textbf{Example:} $\mathrm{Sum}(\cdot)$, $\mathrm{Mean}(\cdot)$ or
    $\mathrm{Max}(\cdot)$ aggregator, e.g.\
    $\vh_v^{(l)}=\mathrm{Sum}(\{\vm_u^{(l)}, u\in N(v)\})$.
  \item \textbf{Issue:} information from node $v$ itself \textbf{could get
    lost} -- computation of $\vh_v^{(l)}$ does not directly depend on
    $\vh_v^{(l-1)}$.
  \item \textbf{Solution:} include $\vh_v^{(l-1)}$ when computing $\vh_v^{(l)}$:
    \begin{itemize}
      \item \textbf{(1) Message:} compute message from node $v$ itself; usually a
        \textbf{different message computation}: neighbors
        $\vm_u^{(l)}=\mW^{(l)}\vh_u^{(l-1)}$, self
        $\vm_v^{(l)}=\mB^{(l)}\vh_v^{(l-1)}$.
      \item \textbf{(2) Aggregation:} after aggregating from neighbors,
        aggregate the message from node $v$ itself, via
        \textbf{concatenation} or \textbf{summation}:
    \end{itemize}
\end{itemize}
\begin{equation}\label{eq:lec04-self}
  \vh_v^{(l)} = \CONCAT\Big(\underbrace{\AGG\big(\{\vm_u^{(l)}, u\in N(v)\}\big)}_{\text{first aggregate from neighbors}},\;
  \underbrace{\vm_v^{(l)}}_{\text{then from node itself}}\Big)
\end{equation}

\begin{definition} [A single GNN layer]\label{def:lec04-layer} 
\leavevmode\\
  \textbf{(1) Message:} each node computes a message
  \[
    \vm_u^{(l)} = \MSG^{(l)}\big(\vh_u^{(l-1)}\big),\quad u\in\{N(v)\cup v\}
  \]
  \textbf{(2) Aggregation:} aggregate messages from neighbors (and self)
  \[
    \vh_v^{(l)} = \AGG^{(l)}\big(\{\vm_u^{(l)}, u\in N(v)\},\,\vm_v^{(l)}\big)
  \]
  \textbf{Nonlinearity (activation):} adds expressiveness; often written
  $\sigma(\cdot)$, e.g.\ $\ReLU(\cdot)$, $\mathrm{Sigmoid}(\cdot)$, \dots; can be
  added to \textbf{message or aggregation}.
\end{definition}

\subsection{Activation (non-linearity)}\label{sec:lec04-act}
\slides{20}

Applied to the $i$-th dimension of embedding $\vx$:
\begin{align}
  \ReLU(\vx_i) &= \max(\vx_i, 0) && \text{most commonly used}\\
  \sigma(\vx_i) &= \frac{1}{1+e^{-\vx_i}} && \text{only to restrict the range of embeddings}\\
  \mathrm{PReLU}(\vx_i) &= \max(\vx_i,0) + a_i\min(\vx_i,0) && \text{$a_i$ trainable}
\end{align}
\begin{itemize}
  \item Plots: ReLU $=0$ for $x<0$, $y=x$ for $x>0$; sigmoid in $(0,1)$;
    PReLU $y=ax$ for $x<0$, $y=x$ for $x>0$.
  \item \textbf{PReLU empirically performs better than ReLU.}
\end{itemize}

\subsection{Classical GNN layers: GCN}\label{sec:lec04-gcn}
\slides{21--22}

\textbf{Graph Convolutional Networks} (Kipf \& Welling, ICLR 2017):
\begin{equation}\label{eq:lec04-gcn}
  \vh_v^{(l)} = \sigma\Big(\mW^{(l)}\sum_{u\in N(v)}\frac{\vh_u^{(l-1)}}{|N(v)|}\Big)
  = \sigma\Big(\underbrace{\sum_{u\in N(v)}}_{\text{Aggregation}}
    \underbrace{\mW^{(l)}\frac{\vh_u^{(l-1)}}{|N(v)|}}_{\text{Message}}\Big)
\end{equation}
\begin{itemize}
  \item \textbf{Message:} each neighbor
    $\vm_u^{(l)}=\frac{1}{|N(v)|}\mW^{(l)}\vh_u^{(l-1)}$ -- \textbf{normalized by
    node degree} (the GCN paper uses a slightly different normalization).
  \item \textbf{Aggregation:} \textbf{sum} over messages from neighbors, then
    apply activation:
    $\vh_v^{(l)}=\sigma\big(\mathrm{Sum}(\{\vm_u^{(l)}, u\in N(v)\})\big)$.
  \item In GCN the input graph is assumed to have \textbf{self-edges} that are
    included in the summation.
\end{itemize}

\subsection{Classical GNN layers: GraphSAGE}\label{sec:lec04-sage}
\slides{23--25}

\textbf{GraphSAGE} (Hamilton et al., NeurIPS 2017):
\begin{equation}\label{eq:lec04-sage}
  \vh_v^{(l)} = \sigma\Big(\mW^{(l)}\cdot\CONCAT\big(\vh_v^{(l-1)},
    \AGG\big(\{\vh_u^{(l-1)}, \forall u\in N(v)\}\big)\big)\Big)
\end{equation}
\begin{itemize}
  \item \textbf{Two-stage aggregation:}
    \begin{itemize}
      \item Stage 1 (aggregate from node neighbors):
        $\vh_{N(v)}^{(l)} \leftarrow \AGG\big(\{\vh_u^{(l-1)}, \forall u\in N(v)\}\big)$
      \item Stage 2 (further aggregate over the node itself):\\
        $\vh_v^{(l)} \leftarrow \sigma\big(\mW^{(l)}\cdot\CONCAT(\vh_v^{(l-1)},\vh_{N(v)}^{(l)})\big)$
    \end{itemize}
  \item \textbf{Message} is computed within the $\AGG(\cdot)$. Neighbor
    aggregation options:
\end{itemize}
\begin{align}
  \text{\textbf{Mean}}:\quad & \AGG = \sum_{u\in N(v)} \frac{\vh_u^{(l-1)}}{|N(v)|}\\
  \text{\textbf{Pool}}:\quad & \AGG = \mathrm{Mean}\big(\{\mathrm{MLP}(\vh_u^{(l-1)}), \forall u\in N(v)\}\big)\\
  \text{\textbf{LSTM}}:\quad & \AGG = \mathrm{LSTM}\big([\vh_u^{(l-1)}, \forall u\in \pi(N(v))]\big)
\end{align}
\begin{itemize}
  \item Mean: weighted average of neighbors (aggregation $\sum$, message
    $\vh_u^{(l-1)}/|N(v)|$). Pool: transform neighbor vectors (message MLP),
    then apply a symmetric vector function $\mathrm{Mean}(\cdot)$ or
    $\mathrm{Max}(\cdot)$ (aggregation). LSTM: apply LSTM (aggregation) to a
    reshuffling $\pi$ of the neighbors.
  \item \textbf{$\ell_2$ normalization} -- \textbf{optional:} apply to
    $\vh_v^{(l)}$ at every layer:
\end{itemize}
\begin{equation}\label{eq:lec04-l2}
  \vh_v^{(l)} \leftarrow \frac{\vh_v^{(l)}}{\norm{\vh_v^{(l)}}_2}\quad\forall v\in V,
  \qquad \norm{u}_2=\sqrt{\textstyle\sum_i u_i^2}\ \ (\ell_2\text{-norm})
\end{equation}
\begin{itemize}
  \item[] Without it, embedding vectors have different scales ($\ell_2$-norm);
    after it, all vectors have the same $\ell_2$-norm. In some cases (not
    always) it results in a performance improvement.
\end{itemize}

\subsection{Classical GNN layers: GAT}\label{sec:lec04-gat}
\slides{26--35}

\textbf{Graph Attention Networks} (Velickovic et al., ICLR 2018; attention:
Vaswani et al., NIPS 2017), with \term{attention weights} $\alpha_{vu}$:
\begin{equation}\label{eq:lec04-gat}
  \vh_v^{(l)} = \sigma\Big(\sum_{u\in N(v)}\alpha_{vu}\mW^{(l)}\vh_u^{(l-1)}\Big)
  = \sigma\Big(\sum_{u\in N(v)}\alpha_{vu}\vm_u^{(l)}\Big)
\end{equation}
\begin{itemize}
  \item \textbf{In GCN / GraphSAGE:} $\alpha_{vu}=\frac{1}{|N(v)|}$ is the
    \textbf{weighting factor (importance)} of node $u$'s message to node $v$
    $\Rightarrow$ $\alpha_{vu}$ is defined \textbf{explicitly} from structural
    properties of the graph (node degree) $\Rightarrow$ all neighbors $u\in
    N(v)$ are equally important to $v$.
  \item \textbf{Not all node's neighbors are equally important.}
    \term{Attention} is inspired by cognitive attention: $\alpha_{vu}$ focuses on
    the important parts of the input data and fades out the rest.
    \begin{itemize}
      \item Idea: the NN should devote more computing power to that small but
        important part of the data; which part is important depends on the
        context and is learned through training.
    \end{itemize}
  \item \textbf{Can we do better than simple neighborhood aggregation? Can the
    weighting factors $\alpha_{vu}$ be learned?}
    \begin{itemize}
      \item \textbf{Goal:} specify \textbf{arbitrary importance} to different
        neighbors of each node in the graph.
      \item \textbf{Idea:} compute $\vh_v^{(l)}$ of each node following an
        \textbf{attention strategy}: nodes attend over their neighborhoods'
        message, implicitly specifying different weights to different nodes in
        a neighborhood.
    \end{itemize}
\end{itemize}

\paragraph{Attention mechanism $a$.} $\alpha_{vu}$ is computed as a byproduct
of an \term{attention mechanism} $a$:
\begin{itemize}
  \item (1) $a$ computes \term{attention coefficients} $e_{vu}$ across pairs of
    nodes $u,v$ based on their messages; \textbf{$e_{vu}$ indicates the
    importance of $u$'s message to node $v$}:
\end{itemize}
\begin{equation}\label{eq:lec04-gat-e}
  e_{vu} = a\big(\mW^{(l)}\vh_v^{(l-1)}, \mW^{(l)}\vh_u^{(l-1)}\big)
         = a\big(\vm_v^{(l)}, \vm_u^{(l)}\big)
\end{equation}
\begin{itemize}
  \item[] Example: $e_{AB} = a\big(\mW^{(l)}\vh_A^{(l-1)}, \mW^{(l)}\vh_B^{(l-1)}\big)
    = a(\vm_A^{(l)},\vm_B^{(l)})$.
\end{itemize}
\begin{itemize}
  \item (2) \textbf{Normalize} $e_{vu}$ into the \textbf{final attention weight}
    $\alpha_{vu}$ with the \textbf{softmax}, so that $\sum_{u\in N(v)}\alpha_{vu}=1$:
\end{itemize}
\begin{equation}\label{eq:lec04-gat-alpha}
  \alpha_{vu} = \frac{\exp(e_{vu})}{\sum_{k\in N(v)}\exp(e_{vk})}
\end{equation}
\begin{itemize}
  \item (3) \textbf{Weighted sum} based on $\alpha_{vu}$:
    $\vh_v^{(l)} = \sigma\big(\sum_{u\in N(v)}\alpha_{vu}\mW^{(l)}\vh_u^{(l-1)}\big)$.
    Example (node $A$, neighbors $B,C,D$):
    \[
      \vh_A^{(l)} = \sigma\big(\alpha_{AB}\mW^{(l)}\vh_B^{(l-1)}
        + \alpha_{AC}\mW^{(l)}\vh_C^{(l-1)} + \alpha_{AD}\mW^{(l)}\vh_D^{(l-1)}\big)
    \]
  \item \textbf{Form of $a$?} The approach is \textbf{agnostic to the choice of
    $a$}; e.g.\ a simple single-layer neural network ($a$ has trainable
    parameters = weights in the Linear layer):
\end{itemize}
\begin{equation}\label{eq:lec04-gat-linear}
  e_{AB} = a\big(\mW^{(l)}\vh_A^{(l-1)}, \mW^{(l)}\vh_B^{(l-1)}\big)
         = \mathrm{Linear}\Big(\mathrm{Concat}\big(\mW^{(l)}\vh_A^{(l-1)}, \mW^{(l)}\vh_B^{(l-1)}\big)\Big)
\end{equation}
\begin{itemize}
  \item \textbf{Parameters of $a$ are trained jointly} with the weight matrices
    (other parameters of the NN, $\mW^{(l)}$) in an end-to-end fashion.
  \item \term{Multi-head attention}: \textbf{stabilizes the learning process} of
    the attention mechanism. Create \textbf{multiple attention scores} (each
    replica with a different set of parameters):
\end{itemize}
\begin{equation}\label{eq:lec04-multihead}
  \vh_v^{(l)}[k] = \sigma\Big(\sum_{u\in N(v)}\alpha_{vu}^{k}\mW^{(l)}\vh_u^{(l-1)}\Big),\ k=1,2,3;
  \qquad
  \vh_v^{(l)} = \AGG\big(\vh_v^{(l)}[1], \vh_v^{(l)}[2], \vh_v^{(l)}[3]\big)
\end{equation}
\begin{itemize}
  \item[] Outputs are aggregated by \textbf{concatenation or summation}.
\end{itemize}
\begin{caution}
  Slide~34 writes the same $\mW^{(l)}$ in every head; in the GAT paper each head
  $k$ also has its own weight matrix $\mW^{k}$.
\end{caution}

\paragraph{Benefits of attention.}
\begin{itemize}
  \item \textbf{Key benefit:} allows for (implicitly) specifying
    \textbf{different importance values ($\alpha_{vu}$) to different neighbors}.
  \item \textbf{Computationally efficient:} attentional coefficients can be
    computed in parallel across all edges; aggregation parallelized across all
    nodes.
  \item \textbf{Storage efficient:} sparse matrix operations need no more than
    $\bigO(V+E)$ entries; \textbf{fixed} number of parameters, irrespective of
    graph size.
  \item \textbf{Localized:} only \textbf{attends over local network
    neighborhoods}.
  \item \textbf{Inductive capability:} a shared \emph{edge-wise} mechanism; does
    not depend on the global graph structure.
\end{itemize}

\begin{table}[htbp]
  \centering\small
  \caption{The classical layers in message + aggregation form (slides~21--34).}
  \label{tab:lec04-layers}
  \renewcommand{\arraystretch}{1.25}
  \begin{tabularx}{\textwidth}{@{}l>{\raggedright\arraybackslash}p{3.9cm}>{\raggedright\arraybackslash}X>{\raggedright\arraybackslash}p{2.9cm}@{}}
    \toprule
    Layer & Message & Aggregation & Weight $\alpha_{vu}$\\
    \midrule
    GCN & $\vm_u^{(l)}=\frac{1}{|N(v)|}\mW^{(l)}\vh_u^{(l-1)}$
        & $\sigma\big(\mathrm{Sum}(\{\vm_u^{(l)}, u\in N(v)\})\big)$ (self-edges
          included)
        & $1/|N(v)|$, fixed (degree)\\
    GraphSAGE & inside $\AGG$: $\vh_u^{(l-1)}/|N(v)|$ (mean),
          $\mathrm{MLP}(\vh_u^{(l-1)})$ (pool)
        & $\sum$/Mean/LSTM over $N(v)$, then
          $\sigma(\mW^{(l)}\cdot\CONCAT(\vh_v^{(l-1)}, \vh_{N(v)}^{(l)}))$
        & $1/|N(v)|$ (mean agg.)\\
    GAT & $\vm_u^{(l)}=\mW^{(l)}\vh_u^{(l-1)}$
        & $\sigma\big(\sum_{u\in N(v)}\alpha_{vu}\vm_u^{(l)}\big)$; heads
          concatenated or summed
        & learned: softmax of $e_{vu}$\\
    \bottomrule
  \end{tabularx}
\end{table}

% ---------------------------------------------------------------------
\section{GNN layers in practice}\label{sec:lec04-practice}
\slides{37--42}

\begin{itemize}
  \item In practice, these classic GNN layers are \textbf{a great starting
    point}; we can often get better performance by \textbf{considering a
    general GNN layer design}, i.e.\ including \textbf{modern deep learning
    modules} that proved useful in many domains (\cref{fig:lec04-suggested}):
    \begin{itemize}
      \item \textbf{Batch normalization:} stabilize neural network training.
      \item \textbf{Dropout:} prevent overfitting.
      \item \textbf{Attention/Gating:} control the importance of a message.
      \item \textbf{More:} any other useful deep learning modules.
    \end{itemize}
\end{itemize}

\begin{figure}[htbp]
  \centering
  \begin{tikzpicture}[lay/.style={draw=BrickRed, thick, text=BrickRed,
      font=\scriptsize\sffamily, minimum height=5.5mm, inner sep=3pt,
      text height=1.6ex, text depth=0.4ex}]
    \node[lay] (l1) at (0,0) {Linear};
    \node[lay, right=4mm of l1] (l2) {BatchNorm};
    \node[lay, right=4mm of l2] (l3) {Dropout};
    \node[lay, right=4mm of l3] (l4) {Activation};
    \node[lay, right=4mm of l4] (l5) {Attention};
    \node[lay, right=4mm of l5] (l6) {Aggregation};
    \draw[arrow] ($(l1.west)+(-0.6,0)$) -- (l1);
    \foreach \a/\b in {l1/l2, l2/l3, l3/l4, l4/l5, l5/l6} \draw[arrow] (\a) -- (\b);
    \draw[arrow] (l6.east) -- ++(0.6,0);
    \draw[decorate, decoration={brace, mirror, amplitude=4pt}]
      ($(l1.south west)+(0,-0.08)$) -- ($(l5.south east)+(0,-0.08)$)
      node[midway, below=5pt, font=\scriptsize\sffamily\bfseries] {Transformation};
  \end{tikzpicture}
  \caption{A suggested GNN layer (slides~37--38, 42).}
  \label{fig:lec04-suggested}
\end{figure}

\paragraph{Batch normalization} (Ioffe \& Szegedy, ICML 2015).
\begin{itemize}
  \item \textbf{Goal:} stabilize neural network training. \textbf{Idea:} given a
    batch of inputs (node embeddings), re-center them to zero mean and re-scale
    the variance to unit variance.
  \item \textbf{Input:} $\mX\in\R^{N\times D}$ ($N$ node embeddings);
    \textbf{trainable parameters:} $\bm{\gamma},\bm{\beta}\in\R^{D}$;
    \textbf{output:} $\mY\in\R^{N\times D}$ (normalized node embeddings).
  \item \textbf{Step 1:} compute the mean and variance over $N$ embeddings;
    \textbf{Step 2:} normalize the feature using the computed mean and variance:
\end{itemize}
\begin{align}
  \bm{\mu}_j &= \frac{1}{N}\sum_{i=1}^{N}\mX_{i,j}, &
  \bm{\sigma}_j^2 &= \frac{1}{N}\sum_{i=1}^{N}\big(\mX_{i,j}-\bm{\mu}_j\big)^2
    \label{eq:lec04-bn1}\\
  \widehat{\mX}_{i,j} &= \frac{\mX_{i,j}-\bm{\mu}_j}{\sqrt{\bm{\sigma}_j^2+\epsilon}}, &
  \mY_{i,j} &= \bm{\gamma}_j\widehat{\mX}_{i,j} + \bm{\beta}_j
    \label{eq:lec04-bn2}
\end{align}

\paragraph{Dropout} (Srivastava et al., JMLR 2014).
\begin{itemize}
  \item \textbf{Goal:} regularize a neural net to prevent overfitting.
    \textbf{During training:} with some probability $p$, randomly set neurons to
    zero (turn off). \textbf{During testing:} use all the neurons.
  \item \textbf{In GNNs}, dropout is applied to \textbf{the linear layer in the
    message function}, e.g.\ $\vm_u^{(l)}=\mW^{(l)}\vh_u^{(l-1)}$: neurons of the
    layer $\vh_u^{(l-1)}\to\vm_u^{(l)}$ are removed.
\end{itemize}

\paragraph{Summary.} Modern deep learning modules can be included into a GNN
layer for better performance. \textbf{Designing novel GNN layers is still an
active research frontier!} Resource: explore diverse GNN designs or try your own
ideas in \textbf{GraphGym}.

% ---------------------------------------------------------------------
\section{Stacking layers of a GNN}\label{sec:lec04-stack}

\subsection{Stacking layers and over-smoothing}\label{sec:lec04-oversmooth}
\slides{44--49}

\begin{itemize}
  \item \textbf{How to connect GNN layers into a GNN?} (3) layer connectivity:
    stack layers sequentially; ways of adding skip connections.
  \item \textbf{Standard way:} stack GNN layers sequentially. \textbf{Input:}
    initial raw node feature $\vx_v$; \textbf{output:} node embeddings
    $\vh_v^{(L)}$ after $L$ GNN layers:
    \[
      \vh_v^{(0)}=\vx_v \xrightarrow{\text{GNN layer}} \vh_v^{(1)}
      \xrightarrow{\text{GNN layer}} \vh_v^{(2)} \xrightarrow{\text{GNN layer}}
      \vh_v^{(3)}\ \cdots
    \]
  \item \textbf{Issue of stacking many GNN layers:} GNNs suffer from
    \term{the over-smoothing problem}: \textbf{all the node embeddings converge
    to the same value}. Bad, because we \textbf{want to use node embeddings to
    differentiate nodes}.
  \item \textbf{Why does over-smoothing happen?}
    \term{Receptive field}: the set of nodes that determine the embedding of a
    node of interest. \textbf{In a $K$-layer GNN, each node has a receptive
    field of $K$-hop neighborhood} (figure: receptive fields for 1-, 2-, 3-layer
    GNNs grow fast on a sample graph).
  \item \textbf{Receptive field overlap} for two nodes: \textbf{the shared
    neighbors quickly grow} when we increase the number of hops (number of GNN
    layers): 1-hop overlap \textbf{only 1 node}; 2-hop \textbf{about 20 nodes};
    3-hop \textbf{almost all the nodes!}
\end{itemize}
\begin{keyidea}[Receptive field \& over-smoothing]
  The embedding of a node is determined by its receptive field; if two nodes have
  highly-overlapped receptive fields, their embeddings are highly similar.\\
  Stack many GNN layers $\to$ nodes have highly-overlapped receptive fields
  $\to$ node embeddings highly similar $\to$ over-smoothing.
\end{keyidea}

\subsection{Lesson 1: be cautious when adding GNN layers}\label{sec:lec04-lesson1}
\slides{50--52}

\begin{itemize}
  \item Unlike neural networks in other domains (CNN for image classification),
    \textbf{adding more GNN layers does not always help}.
    \begin{itemize}
      \item \textbf{Step 1: analyze the necessary receptive field} to solve your
        problem, e.g.\ by computing the diameter of the graph.
      \item \textbf{Step 2:} set the number of GNN layers $L$ to be a bit more
        than the receptive field we like. \textbf{Do not set $L$ to be
        unnecessarily large!}
    \end{itemize}
  \item \textbf{Question:} how to make a GNN more expressive \textbf{if the
    number of GNN layers is small}?
  \item \textbf{Solution 1: increase the expressive power within each GNN
    layer.} In the previous examples each transformation or aggregation function
    only includes one linear layer $\Rightarrow$ \textbf{make aggregation /
    transformation a deep neural network} (if needed, each box -- (1)
    transformation, (2) aggregation -- could include a \textbf{3-layer MLP}).
  \item \textbf{Solution 2: add layers that do not pass messages.} A GNN does
    not necessarily only contain GNN layers: add \textbf{MLP layers} (applied to
    each node) before and after the GNN layers, as \term{pre-process layers} and
    \term{post-process layers} (\cref{fig:lec04-connect}a).
    \begin{itemize}
      \item \textbf{Pre-processing layers:} important when encoding node
        features is necessary, e.g.\ when nodes represent images/text.
      \item \textbf{Post-processing layers:} important when reasoning /
        transformation over node embeddings are needed, e.g.\ graph
        classification, knowledge graphs.
      \item \textbf{In practice, adding these layers works great!}
    \end{itemize}
\end{itemize}

\subsection{Lesson 2: add skip connections}\label{sec:lec04-skip}
\slides{53--56}

\begin{itemize}
  \item What if my problem still requires many GNN layers? $\Rightarrow$
    \textbf{add skip connections in GNNs} (He et al., ResNet, CVPR 2015).
    \begin{itemize}
      \item \textbf{Observation from over-smoothing:} node embeddings in earlier
        GNN layers can sometimes better differentiate nodes.
      \item \textbf{Solution:} increase the impact of earlier layers on the final
        node embeddings, \textbf{by adding shortcuts in the GNN}.
      \item Idea: duplicate $\vx$ into two branches, sum the two branches:
        before adding shortcuts $F(\vx)$; after adding shortcuts
        $F(\vx)+\vx$ (\cref{fig:lec04-connect}b).
    \end{itemize}
  \item \textbf{Why do skip connections work?} Intuition: skip connections
    create \textbf{a mixture of models} (Veit et al., 2016).
    \begin{itemize}
      \item $N$ skip connections $\to$ $2^N$ possible paths (each module: path 1
        includes it, path 2 skips it); each path could have up to $N$ modules.
      \item Example, 3-block residual network: all the possible paths
        $2\cdot2\cdot2=2^3=8$.
      \item We automatically get \textbf{a mixture of shallow GNNs and deep
        GNNs}.
    \end{itemize}
  \item \textbf{Example: GCN with skip connections.}
\end{itemize}
\begin{align}
  \text{standard GCN layer:}\quad
    & \vh_v^{(l)} = \sigma\Big(\underbrace{\textstyle\sum_{u\in N(v)}\mW^{(l)}\frac{\vh_u^{(l-1)}}{|N(v)|}}_{F(\vx)}\Big)
    \label{eq:lec04-gcn-std}\\
  \text{GCN layer with skip connection:}\quad
    & \vh_v^{(l)} = \sigma\Big(\underbrace{\textstyle\sum_{u\in N(v)}\mW^{(l)}\frac{\vh_u^{(l-1)}}{|N(v)|}}_{F(\vx)}
      + \underbrace{\vh_v^{(l-1)}}_{\vx}\Big)
    \label{eq:lec04-gcn-skip}
\end{align}
\begin{itemize}
  \item \textbf{Other options} (Xu et al., jumping knowledge networks, ICML
    2018): \textbf{directly skip to the last layer} -- the final layer directly
    \textbf{aggregates from all the node embeddings} in the previous layers
    (layer aggregation: Concat/Pooling/LSTM), \cref{fig:lec04-connect}c.
\end{itemize}

\begin{figure}[htbp]
  \centering
  \begin{subfigure}[b]{0.34\textwidth}
    \centering
    \begin{tikzpicture}[scale=0.9]
      \foreach \t/\y/\n in {MLP Layer/0/m1, MLP Layer/-0.75/m2, GNN Layer/-1.6/g1,
          GNN Layer/-2.35/g2, GNN Layer/-3.1/g3, MLP Layer/-3.95/m3,
          MLP Layer/-4.7/m4}
        \node[lec04lay] (\n) at (0,\y) {\t};
      \draw[arrow] (0,0.55) -- (m1);
      \foreach \a/\b in {m1/m2, m2/g1, g1/g2, g2/g3, g3/m3, m3/m4} \draw[arrow] (\a) -- (\b);
      \draw[arrow] (m4) -- (0,-5.25);
      \foreach \y in {-1.6,-2.35,-3.1}
        \draw[BrickRed, densely dotted, thick, -{Stealth[length=1.6mm]}]
          (0,\y+0.33) -- (1.3,\y+0.33) -- (1.3,\y-0.31) -- (0.05,\y-0.31);
      \draw[densely dashed, black!50] (-1.15,0.3) rectangle (2.75,-1.05);
      \draw[densely dashed, black!50] (-1.15,-3.65) rectangle (2.75,-5.0);
      \node[font=\scriptsize\sffamily, text=BrickRed, align=left, anchor=west]
        at (1.35,-0.37) {Pre-\\process};
      \node[font=\scriptsize\sffamily, text=BrickRed, align=left, anchor=west]
        at (1.35,-4.32) {Post-\\process};
      \node[font=\scriptsize\sffamily\bfseries, text=BrickRed, align=left, anchor=west]
        at (1.4,-2.35) {Skip\\conn.};
    \end{tikzpicture}
    \caption{Pre/post-process MLP layers and skip connections (slides~52--53, 55).}
  \end{subfigure}\hfill
  \begin{subfigure}[b]{0.29\textwidth}
    \centering
    \begin{tikzpicture}[scale=0.9]
      \node[font=\small] (x) at (0,0.35) {$\vx$};
      \node[lec04lay, minimum width=17mm] (w1) at (0,-0.8) {weight layer};
      \node[lec04lay, minimum width=17mm] (w2) at (0,-2.0) {weight layer};
      \node[circle, draw, thick, inner sep=0pt, minimum size=3mm] (p) at (0,-3.0) {};
      \draw (p.north) -- (p.south) (p.west) -- (p.east);
      \draw[arrow] (0,0.1) -- (w1);
      \draw[arrow] (w1) -- node[right, font=\scriptsize\sffamily] {relu} (w2);
      \draw[arrow] (w2) -- (p);
      \draw[arrow] (p) -- node[right, font=\scriptsize\sffamily] {relu} (0,-3.75);
      \draw[arrow] (0,-0.15) to[out=0, in=0, looseness=1.6]
        node[right, font=\scriptsize, align=center] {$\vx$\\identity} (p.east);
      \node[font=\scriptsize, anchor=east] at (-1.0,-1.4) {$F(\vx)$};
      \node[font=\scriptsize, anchor=east] at (-0.25,-3.0) {$F(\vx)+\vx$};
    \end{tikzpicture}
    \caption{Residual block (slide~53).}
  \end{subfigure}\hfill
  \begin{subfigure}[b]{0.34\textwidth}
    \centering
    \begin{tikzpicture}[scale=0.9]
      \node[font=\scriptsize] (in) at (0,0.25) {Input: $\vh_v^{(0)}$};
      \node[lec04lay] (g1) at (0,-0.75) {GNN Layer};
      \node[lec04lay] (g2) at (0,-1.85) {GNN Layer};
      \node[lec04lay] (g3) at (0,-2.95) {GNN Layer};
      \node[lec04lay, minimum width=30mm, align=center] (agg) at (0,-4.35)
        {Layer aggregation\\Concat/Pooling/LSTM};
      \draw[arrow] (in) -- (g1);
      \draw[arrow] (g1) -- node[left, font=\scriptsize] {$\vh_v^{(1)}$} (g2);
      \draw[arrow] (g2) -- node[left, font=\scriptsize] {$\vh_v^{(2)}$} (g3);
      \draw[arrow] (g3) -- node[left, font=\scriptsize] {$\vh_v^{(3)}$} (agg);
      \draw[arrow] (agg) -- (0,-5.2) node[below, font=\scriptsize] {Output: $\vh_v^{(\mathit{final})}$};
      \foreach \y/\x in {-0.3/1.65, -1.3/1.45, -2.4/1.25}
        \draw[BrickRed, thick, -{Stealth[length=1.6mm]}]
          (0,\y) -- (\x,\y) -- (\x,\y |- agg.north);
    \end{tikzpicture}
    \caption{Jumping to the last layer: aggregate all layers' embeddings (slide~56).}
  \end{subfigure}
  \caption{Layer connectivity (slides~52--56).}
  \label{fig:lec04-connect}
\end{figure}

% ---------------------------------------------------------------------
\bigskip\phantomsection
\addcontentsline{toc}{section}{Lecture 4 at a glance}
\begin{summarybox}[Lecture 4 at a glance]
  \begin{itemize}
    \item \textbf{GNN framework:} (1) message, (2) aggregation, (3) layer
      connectivity, (4) graph augmentation, (5) learning objective.
    \item \textbf{GNN layer:} $\vm_u^{(l)}=\MSG^{(l)}(\vh_u^{(l-1)})$,
      $u\in\{N(v)\cup v\}$; $\vh_v^{(l)}=\AGG^{(l)}(\{\vm_u^{(l)},u\in N(v)\},\vm_v^{(l)})$;
      activation $\sigma$ (ReLU, sigmoid, PReLU).
    \item \textbf{GCN:} $\vh_v^{(l)}=\sigma\big(\sum_{u\in N(v)}\mW^{(l)}\vh_u^{(l-1)}/|N(v)|\big)$.
      \textbf{GraphSAGE:} $\vh_v^{(l)}=\sigma\big(\mW^{(l)}\cdot\CONCAT(\vh_v^{(l-1)},
      \AGG(\{\vh_u^{(l-1)}\}))\big)$, AGG = mean / pool / LSTM, optional $\ell_2$ norm.
    \item \textbf{GAT:} $\vh_v^{(l)}=\sigma\big(\sum_{u\in N(v)}\alpha_{vu}\mW^{(l)}\vh_u^{(l-1)}\big)$,
      $e_{vu}=a(\mW^{(l)}\vh_v^{(l-1)},\mW^{(l)}\vh_u^{(l-1)})$,
      $\alpha_{vu}=\exp(e_{vu})/\sum_{k\in N(v)}\exp(e_{vk})$; multi-head attention.
    \item \textbf{In practice:} Linear $\to$ BatchNorm $\to$ Dropout $\to$
      Activation $\to$ Attention $\to$ Aggregation.
    \item \textbf{Over-smoothing:} $K$ layers $=$ $K$-hop receptive field;
      overlapping fields $\Rightarrow$ similar embeddings. Lesson 1: choose $L$
      just above the needed receptive field; more expressive layers (MLPs),
      pre/post-processing layers. Lesson 2: skip connections $F(\vx)+\vx$
      ($2^N$ paths), skip to the last layer.
  \end{itemize}
\end{summarybox}
